\documentclass[10pt,a4paper]{article}
\usepackage{amsmath,amssymb,amsthm,graphicx}
% Minimal retained preamble for S3 SIGMA article cards taken from
% sources/sigma-2023-029-total-mean-curvature-and-first-dirac-eigenva/pinned-original-source.tex
% (doi:10.3842/SIGMA.2023.029, SIGMA 19 (2023), CC BY 4.0).
%
% Omitted and why: the publisher \documentclass{sigma} block and the
% frontmatter commands - the sigma class is unavailable in the toolchain
% and the assembler supplies the article class.
% The theorem declarations, equation numbering, \allowdisplaybreaks and
% the author's operator macros below are copied verbatim from the pinned
% source preamble (source lines 9-16, 47 and 42-44); of the author's
% macro block at source lines 18-44 only the three macros that appear in
% the delivered spans are retained.
% hyperref is added beyond the source preamble because the author's
% bibliography entries use \href and the unavailable sigma class is
% what provides it in the published version.
\usepackage{amsmath,amssymb,amsthm}
\usepackage{mathtools}
\usepackage{hyperref}
\allowdisplaybreaks
\numberwithin{equation}{section}
\newtheorem{Theorem}{Theorem}[section]
\newtheorem{Corollary}[Theorem]{Corollary}

\theoremstyle{definition}
\newtheorem{Remark}[Theorem]{Remark}

\def\Paps{{\bf P}_+}
\def\Ss{R\!\!\!\!/\,}
\def\D{D\!\!\!\!/\,}
\title{Euclidean Dirac upper bound for mean-convex star-shaped hypersurfaces (Theorem 1.1)}
\author{Raulot, Simon}
\date{doi:10.3842/SIGMA.2023.029; CC BY 4.0}
\begin{document}
\maketitle
\setcounter{section}{1}
\setcounter{Theorem}{0}
\noindent\textit{Editorial note.} All delivered mathematical text is the author's own, copied from the pinned source file sources/sigma-2023-029-total-mean-curvature-and-first-dirac-eigenva/pinned-original-source.tex; no statement or proof step was rewritten, reordered or removed.\par\medskip
\noindent\textit{Editorial note.} The counted span is the author Theorem environment NewUpperBound with its own label and its own single proof environment (source lines 224-258).\par\medskip
\noindent\textit{Editorial note.} The quasi-spherical display (QS-metric) and the mass inequality (GeneralInequality) sit before the statement; the Herzlich theorem (HerzlichPMT-n, printed 2.1) and the corner remark (SpinBoundaryCorner, printed 2.4) are theorem-like contexts placed after the statement so the card keeps 1.1 (both print locally renumbered).\par\medskip
\noindent\textit{Editorial note.} The bibitems EichmairMiaoWang, Herzlich1, Herzlich2 and BartnikChrusciel are extracted verbatim from the pinned source's own thebibliography.\par\medskip
\noindent\textit{Editorial note.} Printed number anchored to the published PDF (SIGMA 19 (2023) 029, read from the stream-extracted text of sigma23-029.pdf downloaded from sigma-journal.com): the rendered labels read Theorem 1.1., Theorem 1.7. and Corollary 1.8. directly, and the PDF destination enumeration Theorem.1.1-1.8 / 2.1-2.4 / 3.1-3.2 / 4.1 / 5.1-5.3 confirms the shared sectionwise theorem-family counter. The pinned tex has five \textbackslash{}section commands and no \textbackslash{}appendix and its section structure matches the published PDF one-to-one (CONTROL 8.8.109 face).\par\medskip
\noindent\textit{Editorial note.} Equation and theorem numbers inside the card are locally renumbered; every cross-reference is a symbolic \textbackslash{}eqref/\textbackslash{}ref and resolves inside the card.\par\medskip
\noindent\textit{Editorial scope.} Single environment card cut from the SIGMA article doi:10.3842/SIGMA.2023.029 (sigma-2023-029-total-mean-curvature-and-first-dirac-eigenva); statement and proof are the author's own text at source lines 113-119 and 224-258.\par\medskip
\noindent\textit{Allowed standard prerequisites (editorial).}\begin{itemize}
\item Quasi-spherical metric construction (QS-metric). Delivered in context (source lines 197-199): for any positive initial value u\_0 on Sigma there is a positive u making g\_u = u\textasciicircum{}2 drho\textasciicircum{}2 + gamma\_rho scalar flat and asymptotically flat with well-defined ADM mass (Shi-Tam, cited at source line 196).. Load-bearing: the metric g\_u the proof builds with u\_0 = H\_0/(2 lambda\_1) in both cases (source lines 225-229 and 253-257).
\item Shi-Tam mass inequality (GeneralInequality). Delivered in context (source lines 209-211): integral H\_0 >= integral u\_0\textasciicircum{}\{-1\} H\_0 as soon as the ADM mass of (E, g\_u) is nonnegative.. Load-bearing: converts mass nonnegativity into the bound (NewUpperBound1) (source line 229).
\item Positive mass with compact inner boundary (Herzlich, Theorem 2.1). Delivered in context (source lines 215-221): the ADM mass is nonnegative when lambda\_1 >= H/2, and zero mass forces flatness with constant boundary mean curvature.. Load-bearing: nonnegativity at source line 229 and the equality analysis at source lines 229-233 and 257.
\item Eichmair-Miao-Wang deformation. Named construction (cited at source lines 235, 239, 245): a star-shaped hypersurface with positive scalar and mean curvatures is deformed to strictly convex parallels Sigma\_rho with a scalar-flat quasi-spherical metric g\_v on Sigma x [0, rho\_0].. Load-bearing: reduces the star-shaped case to the strictly convex one (source lines 235-257).
\item Positive mass across a Lipschitz corner (SpinBoundaryCorner, Remark 2.4). Delivered in context (source lines 285-297): the positive mass theorem holds for asymptotically flat manifolds with corners along a hypersurface.. Load-bearing: extends the Herzlich theorem to the glued manifold with Lipschitz metric (source line 257).
\end{itemize}\medskip
\begin{gather}\label{QS-metric}
g_u:=u^2\,{\rm d}\rho^2+\gamma_\rho
\end{gather}

\begin{gather}\label{GeneralInequality}
\int_\Sigma \mathcal{H}_0\,{\rm d}\Sigma\geq \int_\Sigma u^{-1}_0\mathcal{H}_0\,{\rm d}\Sigma.
\end{gather}

\begin{Theorem}\label{NewUpperBound}
Let $\Sigma$ be a compact, orientable, mean-convex and star-shaped hypersurface with positive scalar curvature in the $n$-dimensional Euclidean space. Then the first nonnegative eigenvalue $\lambda_1(\Sigma,\gamma)$ of the extrinsic Dirac operator $\D_\gamma$ of $(\Sigma,\gamma)$ satisfies
\begin{gather}\label{NewUpperBound1}
\lambda_1(\Sigma,\gamma)\leq\frac{\int_\Sigma \mathcal{H}_0\,{\rm d}\Sigma}{2|\Sigma|}.
\end{gather}
Here $\mathcal{H}_0$ denotes the mean curvature function of $\Sigma$ in $\mathbb{R}^{n}$ and $|\Sigma|$ its volume. Equality occurs if and only if the hypersurface is a round sphere.
\end{Theorem}

\begin{Theorem}[\cite{Herzlich1,Herzlich2}]\label{HerzlichPMT-n}
Let $(M,g)$ be a $n$-dimensional Riemannian spin asymptotically flat manifold with nonnegative scalar curvature and with a compact inner boundary $\Sigma$. Assume that the first nonnegative eigenvalue $\lambda_1(\Sigma,\gamma)$ is positive and satisfies
\begin{gather}\label{DiracMeanCur-n}
\lambda_1(\Sigma,\gamma)\geq\frac{1}{2}H,
\end{gather}
where $H$ is the mean curvature of $\Sigma$ in $M$. Then the mass is nonnegative and if the mass is zero, $(M,g)$ is flat and the mean curvature is constant and equal to $H=2\lambda_1(\Sigma,\gamma)$.
\end{Theorem}

\begin{Remark}\label{SpinBoundaryCorner}
We briefly explain here how to prove Theorem \ref{HerzlichPMT-n} for asymptotically flat manifolds with corners along a hypersurface. An asymptotically flat manifold $(M,g)$ with boundary~$\Sigma$ is said to have a corner along a hypersurface $N$ if it can be written as the disjoint union of two subsets $M_+$ and $M_-$, where $(M_+,g_{+})$ is a smooth $n$-dimensional asymptotically flat manifold with boundary $N$ and $(M_-,g_-)$ is a smooth $n$-dimensional compact Riemannian manifold with boundary $\Sigma\cup N$ and where $g_\pm=g_{|M_\pm}$. The ADM mass of $(M,g)$ corresponds to the ADM mass of $(M_+,g_+)$. Let also $H_+$ and $H_-$ the mean curvatures of $N$ in $(M_+,g_+)$ and $(M_-,g_-)$ with respect to the unit normal pointing toward infinity. The positive mass theorem for such manifolds asserts that if the scalar curvature of $(M,g)$ is nonnegative and if $H_-\geq H_+$, then its ADM mass is nonnegative. The proof of the positive mass Theorem~\ref{HerzlichPMT-n} relies essentially on two important points: the existence of a harmonic spinor field satisfying the Atiyah--Patodi--Singer condition which is asymptotic to a constant spinor, usually called a Witten spinor, and an integral version of the Schr\"odinger--Lichnerowicz formula. For the first point, it is enough to observe that since the metric $g$ is Lipschitz, the work of Bartnik and Chru\'sciel~\cite{BartnikChrusciel} still applies to construct an adequate spinor $\Phi$. On the other hand, recall that the Schr\"odinger--Lichnerowicz formula asserts that
\begin{gather*}%\label{SL-Formula}
D^2\psi=\nabla^*\nabla\psi+\frac{R}{4}\psi
\end{gather*}
for all spinor field $\psi$ on $M$ and where $D$, $\nabla$ and $\nabla^\ast$ are respectively the Dirac operator, the spin Levi-Civita connection and its $L^2$-formal adjoint on $(M,g)$. Although this formula is no longer valid on $(M,g)$ since the scalar curvature is not defined on $N$, it remains true on $(M_+,g_+)$ and on $(M_-,g_-)$ separately. We can then integrate this formula on $M_-$ first and then on a~domain~$M_R$ in~$M_+$ whose boundary is the union of $N$ and a coordinate sphere at infinity with radius~$R$. Adding them together and inserting the aforementioned Witten spinor $\Phi$ lead to
\begin{align*}
\frac{1}{2}(n-1)\omega_{n-1}m_{\rm ADM}(M,g) ={} & \int_{M}\bigg(|\nabla\Phi|^2+\frac{R}{4}|\Phi|^2\bigg)\,{\rm d}M+\frac{1}{2}\int_N(H_--H_+)|\Phi|^2\,{\rm d}N
\\
& -\int_{\Sigma}\langle \D_\gamma\Phi+\frac{H}{2}\Phi,\Phi\rangle \,{\rm d}\Sigma
\end{align*}
as $R$ goes to infinity. The nonnegativity of the mass then follows from the assumptions $R\geq 0$, $H_-\geq H_+$ and~\eqref{DiracMeanCur-n}.
\end{Remark}

\begin{proof}[Proof of Theorem \ref{NewUpperBound}]
First assumed that $(\Sigma,\gamma)$ is a strictly convex hypersurface in $\mathbb{R}^n$. As~recalled above, one can solve the quasi-spherical metric problem on $\mathcal{E}$ with the initial value
\begin{gather*}%\label{InitialValueBV}
u_0=\frac{1}{2\lambda_1(\Sigma,\gamma)}\mathcal{H}_0>0.
\end{gather*}
Then the metric $g_u$ defined by~\eqref{QS-metric} yields, on $\mathcal{E}$, an asymptotically flat metric with zero scalar curvature for which the inner compact boundary is isometric to $(\Sigma,\gamma)$ and has constant mean curvature $H=2\lambda_1(\Sigma,\gamma)$. It is immediate to see that Theorem \ref{HerzlichPMT-n} applies so that we deduce from~\eqref{GeneralInequality} the inequality~\eqref{NewUpperBound1}. If now we assume that equality occurs, then the ADM mass of $(\mathcal{E},g_u)$ is zero and the positive mass theorem implies that the metric $g_u$ is flat. From the Gauss formula, one can compute that the Riemann curvature tensor ${\rm Riem}_u$ of $g_u$ is given by
\begin{gather*}%\label{RiemFlat}
{\rm Riem}_u(e_i,e_j,e_i,e_j)=\big(1-u^{-2}\big){\rm Riem}_\rho(e_i,e_j,e_i,e_j),
\end{gather*}
where ${\rm Riem}_\rho$ is the Riemann curvature tensor of $\Sigma_\rho$ and $\{e_i\mid 1\leq i\leq n-1\}$ is a local orthonormal frame on $(\Sigma,\gamma)$ parallel translated in the direction of $\frac{\partial}{\partial \rho}$. Since the metric $g_u$ is flat and $\Sigma_\rho$ is strictly convex in~$\mathbb{R}^n$ for all $\rho\geq0$, we conclude that $u\equiv 1$. Then $(\Sigma,\gamma)$ is a smooth embedded compact hypersurface in the Euclidean space with constant mean curvature and so it is a round sphere. The converse is obvious since for a Euclidean round sphere with radius $r>0$, its first Dirac eigenvalue is equal to~$(n-1)/(2r)$ and its mean curvature is equal to $(n-1)/r$.

If now we assume that the hypersurface $\Sigma$ is star-shaped with positive scalar and mean curvatures, we apply the method developed by Eichmair, Miao and Wang~\cite{EichmairMiaoWang}. Under these assumptions, there exists a smooth map $F\colon\Sigma\times[0,\infty)\rightarrow \mathbb{R}^n$ such that
\begin{gather*}
\frac{\partial F}{\partial \rho}=\frac{\mathcal{H}_\rho}{\Ss_\rho}\nu,
\end{gather*}
$F(\Sigma,0)=\Sigma$ and $\Sigma_\rho:=F(\Sigma,\rho)$ has positive mean and scalar curvatures and is a strictly convex hypersurface for $\rho$ sufficiently large, say for $\rho\geq \rho_0$. Here $\nu$ is the outer unit normal to the hypersurface $\Sigma_\rho$, $\mathcal{H}_\rho$ is its mean curvature and $\Ss_\rho$ its scalar curvature. Moreover, the pull-back of the Euclidean metric by $F$ on $\Sigma\times[0,\rho_0]$ has the form $g_\eta$ given by~\eqref{QS-metric} with $\eta$ a smooth positive function. From~\cite[Proposition 2]{EichmairMiaoWang}, there exists a smooth positive function $v$ on $\Sigma\times [0,\rho_0]$ such that the scalar curvature of the metric $g_v$
is zero and with initial data
\begin{gather*}
v(\cdot,0)=\frac{\eta(\cdot,0)}{2\lambda_1(\Sigma,\gamma)}\mathcal{H}_0.
\end{gather*}
This last condition implies that the mean curvature of $\Sigma\simeq\Sigma\times\{0\}$ is equal to $h_0=2\lambda_1(\Sigma,\gamma)$ for the metric $g_v$.
Then, as proved in~\cite[Proposition~3]{EichmairMiaoWang}, the function
\begin{gather*}
\rho\in[0,\rho_0]\mapsto\int_{\Sigma_\rho}\big(\mathcal{H}_\rho-h_\rho\big)\,{\rm d}\Sigma_\rho,
\end{gather*}
where $h_\rho$ represents the mean curvature of $\Sigma_\rho$ with respect to the metric $g_v$, is monotone nonincreasing in $\rho$ and so it follows that
\begin{gather}\label{Monotone1}
\int_{\Sigma}\big(\mathcal{H}_0-h_0\big)\,{\rm d}\Sigma=\int_{\Sigma}\mathcal{H}_0\,{\rm d}\Sigma-2\lambda_1(\Sigma,\gamma)|\Sigma| \geq \int_{\Sigma}(\mathcal{H}_{\rho_0}-h_{\rho_0})\,{\rm d}\Sigma.
\end{gather}
Now we can apply the original construction of Shi and Tam to the strictly convex hypersurface~$\Sigma_{\rho_0}$ by solving the quasi-spherical equation in such a way that the resulting manifold $(\Sigma\times[\rho_0,\infty),g_u)$ is an asymptotically flat manifold whose inner boundary is isometric to $(\Sigma,\gamma_{\rho_0})$ with mean curvature equal to $h_{\rho_0}$. As before, we get from~\eqref{GeneralInequality} that
\begin{gather}\label{Monotone2}
\int_{\Sigma}(\mathcal{H}_{\rho_0}-h_{\rho_0})\,{\rm d}\Sigma\geq c(n)m_{\rm ADM}(\mathcal{E},g_u).
\end{gather}
Gluing $(\Sigma\times[0,\rho_0],g_v)$ and $(\Sigma\times[\rho_0,\infty),g_u)$ along their common boundary $(\Sigma,\gamma_{\rho_0})$ yields an asymptotically flat manifold $(\mathcal{E},\widetilde{g})$ with Lipschitz metric along $\Sigma_{\rho_0}$, with ADM mass equal to $m_{\rm ADM}(\mathcal{E},g_u)$ and with a compact inner boundary whose mean curvature is equal to $2\lambda_1(\Sigma,\gamma)$. Combining~\eqref{Monotone1} and~\eqref{Monotone2} with the fact that Theorem \ref{HerzlichPMT-n} holds in this context (see Remark \ref{SpinBoundaryCorner}) yields the desired inequality. If equality occurs, we deduce as before that the manifold $(\mathcal{E},\widetilde{g})$ is flat away from $\Sigma_{\rho_0}$ and that $u\equiv 1$ since $\Sigma_\rho$ is convex for $\rho\geq \rho_0$. On the other hand, we also deduce from the Gauss formula that $u\equiv \eta$ since $\Sigma_\rho$ has positive scalar curvature for $\rho\in[0,\rho_0]$. This allows to conclude that $(\mathcal{E},\widetilde{g})$ is isometric to the exterior of $\Sigma$ in the Euclidean space. The end of the proof then proceeds as in the previous case.
\end{proof}

\begin{thebibliography}{99}

\bibitem{EichmairMiaoWang}
Eichmair M., Miao P., Wang X., Extension of a theorem of {S}hi and {T}am,
 \href{https://doi.org/10.1007/s00526-011-0402-2}{\textit{Calc. Var. Partial Differential Equations}} \textbf{43} (2012),
 45--56, \href{https://arxiv.org/abs/0911.0377}{arXiv:0911.0377}.


\bibitem{Herzlich1}
Herzlich M., A {P}enrose-like inequality for the mass of {R}iemannian
 asymptotically flat manifolds, \href{https://doi.org/10.1007/s002200050159}{\textit{Comm. Math. Phys.}} \textbf{188}
 (1997), 121--133.


\bibitem{Herzlich2}
Herzlich M., Minimal surfaces, the {D}irac operator and the {P}enrose
 inequality, in S\'eminaire de {T}h\'eorie {S}pectrale et {G}\'eom\'etrie, {A}nn\'ee 2001--2002, \textit{S\'emin. Th\'eor. Spectr. G\'eom.}, Vol.~20,
 Universit\'e de Grenoble I, Saint-Martin-d'H\`eres, 2002, 9--16.


\bibitem{BartnikChrusciel}
Bartnik R.A., Chru\'sciel P.T., Boundary value problems for {D}irac-type
 equations, \href{https://doi.org/10.1515/crll.2005.2005.579.13}{\textit{J.~Reine Angew. Math.}} \textbf{579} (2005), 13--73,
 \href{https://arxiv.org/abs/math.DG/0307278}{arXiv:math.DG/0307278}.


\end{thebibliography}
\end{document}
